\documentclass[instructions.tex]{subfiles}
\begin{document}
 
\chapter{Des veillées, assemblées et promenades nocturnes.}
 
Ce sont là les moyens puissans qu’emploie l’esprit des ténèbres pour tendre efficacement des pièges à l’innocence, et pour pervertir les âmes.

\primo Une gaîté peu mesurée préside d’ordinaire aux veillées, aux assemblées, aux promenades qui se font durant la nuit entre des personnes de différent sexe; il est difficile, et il n’arrive presque jamais que la conversation y soit constamment ce qu’elle doit être. Pour ne point parler de beaucoup d’autres défauts qui s’y rencontrent souvent, on s’y permet presque toujours des discours indiscrets, équivoques, licencieux et capables de produire les plus funestes impressions. A quels dangers ne s’y trouvent point alors exposés de jeunes coeurs encore inexpérimentés et sans défiance? D’abord, on a honte d’entendre de pareils discours; mais peu à peu on s’y habitue; bientôt on y applaudit avec répugnance et par respect humain, mais ensuite par inclination et par goût; enfin, on en tient soi-même, pour se mettre à la hauteur des autres, pour gagner aussi à son tour, fournir du sien et faire de l’esprit. Prétendra-t-on, après cela, sortir de ces assemblées, de ces veillées ou de ces promenades, tel qu’on y était entré, avec un coeur encore innocent? Ce n’est pas ainsi qu’en juge le tentateur lui-même. Il dit de vous avoir conduit dans ces réunions, applaudir au mauvais goût que vous y avez pris, à des plaies profondes que votre vertu y a reçues, au poison mortuif qui y est entré insensiblement dans votre coeur. Il prévoit les suites immanquables où d’autres entraînent un commencement si déplorable; d’après votre expérience longue et souvent répétée, qu’une parole équivoque suffit quelquefois seule pour donner naissance à une foule de pensées, de réflexions, d’affections, de mauvais désirs, qui font le tourment de la vie, et qui deviennent comme une source inépuisable de tentations et de péchés. Déjà il compte en quelque sorte que vous serez un jour associé à son malheur éternel.

Jeunesse chrétienne, trompez l’attente de ce cruel ennemi. Voulez-vous conserver tout de bon le précieux trésor de la sainteté dans les moeurs, évitez à jamais les réunions pleines de dangers dont il s’agit: il est à craindre que si vous vous y rencontrez une seule fois, vous n’en remportiez votre coeur changé, séduit, perverti: tant est vrai ce que dit le grand Apôtre, que les mauvais entretiens corrompent les bonnes moeurs.\footnote{1. Cor. 15. 33.}

\secundo Cependant il est rare qu’on s’en tienne à ces conversations déjà si pernicieuses; souvent on y ajoute des jeux charmans, coupables; souvent on y joint des familiarités indécentes, des embrassemens passionnés. De semblables pratiques réveilleraient les passions mortifiées d’un anachorete qui aurait fait pénitence pendant cinquante ans dans un affreux désert, et elles ne produiraient aucun effet funeste sur des coeurs jeunes, sans expérience, sans précautions! (ce sont là, disent les Pères de la vie spirituelle, les signes d’une chasteté mourante, ou déjà morte;) et on se les permettrait sans scrupule, à un âge où les mesures les plus sévères, la fuite des dangers la plus soigneuse, la vigilance la plus exacte et la prière habituelle sont nécessaires pour assurer la victoire et conserver la délicate vertu! Un regard imprudent est du plus dangereux, et d’un roi qui était prophète, un adultère et un homicide: et des libertés folâtres entre de jeunes personnes de différent sexe seraient pour elles sans danger, sans conséquence! Oser l’espérer, ne serait-ce pas le comble de la présomption? Agir ensuite de cette présomption ne serait-ce pas mériter d’éprouver tout l’effet de cette menace de l’Ecriture: Celui qui aime le danger y périra.\footnote{Eccl. 3. 27.}

\tertio Mais c’est surtout dans les promenades et les veillées, où l’on se trouve seul à seul, que les dangers se multiplient et qu’ils deviennent le plus séduisans. Alors le coeur parle au coeur, les ténèbres inspirent de la sécurité, et l’isolement où l’on se voit, déplorable n’apprend que trop dans quel abîme de péchés et de crimes se laissent entraîner peu à peu, quelquefois même tout d’un coup, des jeunes gens qui fuient ainsi l’oeil vigilant de leurs parens, ou d’autres témoins sages. Faut-il s’étonner si l’on voit cette jeunesse perdue changer sensiblement de caractere, devenir hautaine, arrogante, indocile, d’une humeur âpre, tracassière; si elle perd le goût de la prière, des lectures pieuses, des exercices publics de la religion, des sacremens, de la religion elle-même; et si enfin elle s’associe bientôt aux libertins d’esprit qui rejettent les dogmes sacrés de la foi catholique, puisqu’elle appartient déjà à cette classe aveugle par le libertinage outré du coeur: de là les discordes qui divisent les familles, le déshonneur qui s’y introduit, les scandales qui se répandent au-dehors, les unions mal assorties, illégitimes, contraires aux intérêts de la religion, et dans lesquelles on traîne une vie malheureuse, image vivante de l’enfer vers lequel on se précipite à grands pas.

Jeunes gens, tenez-vous dans le sein de vos familles le soir, ou, si vous en sortez, que ce soit toujours accompagnés de quelques parens sages, prudens, âgés, expérimentés, et qui ne vous perdent pas de vue. Abstenez-vous de toute assemblée où ne régnerait pas la décence la plus exacte dans les paroles, dans les manieres, la tenue, la parure. Si vous pensez à un établissement, ne vous trouvez jamais seul avec la personne que vous avez en vue: quand on ne veut pas faire de mal, on ne craint pas les témoins. Une fille bien élevée, honnête, craignant Dieu, aimant la vertu, ne parlera point à un jeune homme à l’insu de ses parens, et hors de leur présence, ou de quelqu’un qui lui en tienne lieu. Point de promenade le soir, seul à seul: vous vous exposeriez à un grand danger de vous perdre, de perdre aussi votre réputation, et vous donneriez du scandale. Quiconque fait mal hait la lumiere, dit Jésus-Christ dans l’Evangile\footnote{(1) Joan. 3. 20.}. O mon Dieu! répandez sur la jeunesse fidèle l’abondance de vos grâces, afin que, connaissant les dangers qui l’environnent de toutes parts, elle les évite autant qu’elle doit, et que, s’y trouvant engagée sans les avoir cherchés, elle s’en retire promptement, le coeur rempli d’amour pour vous et pour la vertu.

Jusqu’ici nous n’avons parlé que de la jeunesse et qu’à elle seule; qu’on n’en conclue pas que les dangers que nous avons montrés, et que les mesures que nous indiquons ne regardent que cet âge; ce serait s’abuser étrangement. La guerre du chrétien dure toute sa vie; il doit donc se tenir constamment en garde contre les ennemis du salut, se défier continuellement de sa faiblesse, et ne jamais poser les armes, ni s’appuyer sur ses anciens lauriers; l’exemple de Salomon doit effrayer tous les âges et les hommes les plus saints.
 
\section{Résolutions}
 
\primo J’éviterai soigneusement toutes les compagnies dans lesquelles on se permettrait de tenir des discours contraires aux bonnes moeurs.
 
\secundo Quand on en tiendra de cette nature en ma présence, je me hâterai de remplir les devoirs qui me sont imposés dans ce cas: je ferai taire la langue qui s’émancipe, si j’ai de l’autorité, ou quelque influence sur elle, ou si je le puis avec prudence; j’élèverai mon coeur à Dieu; je montrerai par un air sérieux et triste que je ne prends aucune part à ces discours et, s’il m’est possible, je me retirerai incontinent du danger. 

\prayer{Mon Dieu! ne permettez pas que je profane jamais ma langue, en la laissant servir à proférer de mauvaises paroles; ni mes oreilles, en les y ouvrant volontairement; ni mon coeur, en y mettant une affection coupable.}
 
\end{document}
